\documentclass[aspectratio=169,11pt,xcolor=table]{beamer}

\usetheme[progressbar=frametitle, numbering=fraction,
          sectionpage=none, subsectionpage=none, block=fill]{metropolis}
\usepackage[brazilian]{babel}
\usepackage{fontspec}
\setsansfont{Fira Sans}[BoldFont={Fira Sans SemiBold},
  ItalicFont={Fira Sans Italic}, BoldItalicFont={Fira Sans SemiBold Italic}]
\setmonofont{Inconsolatazi4-Regular.otf}[BoldFont=Inconsolatazi4-Bold.otf, Scale=1.08]
\usepackage{listings}
\usepackage{booktabs}
\usepackage{tabularx}
\usepackage{fontawesome5}
\usepackage{appendixnumberbeamer}
\usepackage{tikz}
\usetikzlibrary{positioning,arrows.meta,calc,shapes.geometric}

\definecolor{ink}{HTML}{1E2633}
\definecolor{paper}{HTML}{FBFAF7}
\definecolor{accent}{HTML}{C8553D}
\definecolor{muted}{HTML}{6B7280}
\definecolor{codebg}{HTML}{F2EFE9}
\definecolor{term}{HTML}{1B1F27}
\definecolor{ok}{HTML}{2F7D4F}
\definecolor{bad}{HTML}{B23A3A}
\definecolor{str}{HTML}{3D7A5A}

\setbeamercolor{normal text}{fg=ink, bg=paper}
\setbeamercolor{background canvas}{bg=paper}
\setbeamercolor{frametitle}{fg=paper, bg=ink}
\setbeamercolor{palette primary}{fg=paper, bg=ink}
\setbeamercolor{alerted text}{fg=accent}
\setbeamercolor{progress bar}{fg=accent, bg=ink!25}
\setbeamercolor{title separator}{fg=accent}
\setbeamercolor{block title}{fg=ink, bg=ink!9}
\setbeamercolor{block body}{fg=ink, bg=ink!4}
\setbeamercolor{block title alerted}{fg=paper, bg=accent}
\setbeamercolor{block body alerted}{fg=ink, bg=accent!8}
\setbeamercolor{footnote}{fg=muted}
\setbeamertemplate{itemize item}{\color{accent}\small$\blacktriangleright$}
\setbeamertemplate{itemize subitem}{\color{muted}\small--}

\newcommand{\certo}{\textcolor{ok}{\faCheck}}
\newcommand{\errado}{\textcolor{bad}{\faTimes}}
\newcommand{\fonte}[1]{\par\vfill{\scriptsize\color{muted}#1}}
\newcommand{\destaque}[1]{%
  \begin{center}\colorbox{accent!10}{\parbox{0.92\linewidth}{\centering\strut #1\strut}}\end{center}}
\newcolumntype{Y}{>{\raggedright\arraybackslash}X}
\newcolumntype{P}{>{\columncolor{accent!10}\raggedright\arraybackslash}X}

\lstdefinelanguage{Pl}{
  sensitive=true,
  morecomment=[l]\%,
  morestring=[b]',
  alsoletter={_},
  emph={Id,Conta,Origem,Destino,Valor,Categoria,Limite,Motivo,Alerta,
        A,B,C,D,I,L,X,Y,Lista,Atual,Inicio,Visitados,Prox,Resto,Contas,
        Antes,Achados,PorTransacao,Ciclos,Ids,Regra},
  literate={:-}{{\textbf{:-}}}2 {?-}{{\textbf{?-}}}2,
}
\lstset{
  basicstyle=\ttfamily\footnotesize,
  columns=fullflexible, keepspaces=true, upquote=true,
  showstringspaces=false, extendedchars=true,
  backgroundcolor=\color{codebg},
  frame=single, framerule=0pt, framesep=5pt,
  xleftmargin=5pt, xrightmargin=5pt,
  aboveskip=3pt, belowskip=3pt,
  commentstyle=\color{muted}\itshape,
  stringstyle=\color{str},
  keywordstyle=\color{ink}\bfseries,
  emphstyle=\color{accent},
}
\lstdefinestyle{pl}{language=Pl}
\lstdefinestyle{pequeno}{basicstyle=\ttfamily\scriptsize}
\lstdefinestyle{terminal}{
  backgroundcolor=\color{term},
  basicstyle=\ttfamily\scriptsize\color{paper},
  framesep=7pt,
}

\title{Programação Lógica com Prolog}
\subtitle{Validação de transações financeiras: pensar em fatos e regras}
\author{Grupo 3 · Thiago Martins, Carlos Arthur e Gabriel Marconato Momesso}
\institute{Linguagens de Programação}
\date{2026}

\begin{document}

\maketitle

\begin{frame}{Roteiro}
  \begin{enumerate}
    \setlength{\itemsep}{6pt}
    \item \textbf{O problema e a aplicação}
    \item \textbf{O paradigma lógico}: fatos, regras, unificação, backtracking, recursão
    \item \textbf{Prolog na prática}: como o núcleo de validação foi escrito
    \item \textbf{Comparação} com Java, Python e SQL
    \item \textbf{Dificuldades, onde o paradigma brilha e conclusões}
  \end{enumerate}
  \vspace{8pt}
  \destaque{Pergunta que guia o trabalho: \emph{como o paradigma muda a maneira de
  \textbf{pensar}, \textbf{modelar} e \textbf{implementar} uma solução?}}
\end{frame}

\section{O problema e a aplicação}

\begin{frame}{O problema: validar um lote de transações}
  \textbf{Entrada:} transações \texttt{(id, origem, destino, valor, categoria)},
  limites por categoria e contas de risco.\\
  \textbf{Saída:} \alert{cada problema encontrado, com o motivo}.
  \vspace{8pt}

  \footnotesize
  \begin{tabularx}{\linewidth}{@{}l Y l@{}}
    \toprule
    \textbf{Regra} & \textbf{Viola quando} & \textbf{No lote da proposta} \\
    \midrule
    Limite da categoria & valor acima do limite da categoria & t1, t3 \\
    Conta de risco & valor alto para uma conta de risco & t3 \\
    Autotransferência & origem igual ao destino & t2 \\
    Valor inválido & valor zero ou negativo & t6 \\
    Ciclo de dinheiro & o dinheiro sai de uma conta e volta a ela &
      pedro $\to$ bruno $\to$ carla $\to$ pedro \\
    \bottomrule
  \end{tabularx}
  \normalsize
  \destaque{As regras chegam como frases: \emph{``é suspeito \textbf{se}\ldots''}.
  É exatamente assim que se escreve Prolog.}
\end{frame}

\begin{frame}{A aplicação}
  \centering
  \begin{tikzpicture}[
      box/.style={draw=ink!45, rounded corners=3pt, fill=white, align=center,
                  minimum width=3.0cm, minimum height=1.35cm, font=\small},
      arr/.style={-{Stealth[length=2.2mm]}, ink!55, thick},
      lbl/.style={font=\scriptsize, text=muted}]
    \node[box] (tela) at (0,0)
      {\textbf{Tela}\\[-2pt]{\scriptsize React}\\[-2pt]{\scriptsize\itshape\color{muted} mostra}};
    \node[box] (bff) at (4.6,0)
      {\textbf{BFF}\\[-2pt]{\scriptsize TanStack Start}\\[-2pt]{\scriptsize\itshape\color{muted} busca e envia}};
    \node[box, draw=accent, line width=1.5pt, fill=accent!9] (pl) at (9.8,0)
      {\textbf{Núcleo}\\[-2pt]{\scriptsize SWI-Prolog}\\[-2pt]{\scriptsize\itshape\color{accent} decide}};
    \node[box] (db) at (4.6,-2.2)
      {\textbf{Banco}\\[-2pt]{\scriptsize PostgreSQL}\\[-2pt]{\scriptsize\itshape\color{muted} guarda}};
    \draw[arr] ([yshift=4pt]tela.east) -- node[above, lbl]{pede} ([yshift=4pt]bff.west);
    \draw[arr] ([yshift=-4pt]bff.west) -- node[below, lbl]{achados} ([yshift=-4pt]tela.east);
    \draw[arr] ([yshift=4pt]bff.east) -- node[above, lbl]{POST /validar} ([yshift=4pt]pl.west);
    \draw[arr] ([yshift=-4pt]pl.west) -- node[below, lbl]{achados (JSON)} ([yshift=-4pt]bff.east);
    \draw[arr] ([xshift=-4pt]bff.south) -- node[left, lbl]{SQL} ([xshift=-4pt]db.north);
    \draw[arr] ([xshift=4pt]db.north) -- node[right, lbl]{lotes, limites} ([xshift=4pt]bff.south);
  \end{tikzpicture}

  \vspace{6pt}
  \raggedright\small
  \begin{itemize}
    \item \textbf{Banco:} só a estrutura dos dados, sem \texttt{CHECK} de regra de negócio
    \item \textbf{BFF:} consultas SQL e um \texttt{fetch}; nenhuma comparação de valor ou limite
    \item \textbf{Tela:} destaca apenas os IDs que o núcleo devolve
  \end{itemize}
  \destaque{\textbf{Nenhuma linha fora do Prolog decide se algo é válido.}}
\end{frame}

\section{O paradigma lógico}

\begin{frame}{A pergunta muda}
  \begin{columns}[T]
    \begin{column}{0.47\textwidth}
      \begin{block}{Imperativo / POO}
        \small\textbf{``Como eu calculo a resposta?''}
        \begin{itemize}
          \item passo a passo, em ordem
          \item laços e variáveis que mudam
          \item o programador controla a busca
        \end{itemize}
      \end{block}
    \end{column}
    \begin{column}{0.47\textwidth}
      \begin{alertblock}{Lógico}
        \small\textbf{``O que é verdade sobre o problema?''}
        \begin{itemize}
          \item fatos e regras
          \item variáveis não mudam depois de ligadas
          \item a linguagem procura as respostas
        \end{itemize}
      \end{alertblock}
    \end{column}
  \end{columns}
  \vspace{16pt}
  \begin{center}
    {\Large \textbf{Algoritmo} = \alert{\textbf{Lógica}} + \textbf{Controle}}\\[3pt]
    {\small\color{muted} Robert Kowalski, 1979}\\[8pt]
    {\small nós escrevemos a lógica; o Prolog cuida do controle}
  \end{center}
\end{frame}

\begin{frame}{De onde vem e qual Prolog usamos}
  \centering
  \begin{tikzpicture}[x=2.45cm]
    \draw[ink!35, line width=1.2pt] (-0.3,0) -- (5.3,0);
    \foreach \x/\ano/\txt in {
        0/1965/{Robinson:\\princípio da\\resolução},
        1/1972/{Colmerauer\\e Roussel criam\\o Prolog},
        2/1974/{Kowalski: lógica\\como linguagem\\de programação},
        3/1995/{padrão\\ISO 13211},
        4/2011/{IBM Watson\\usa Prolog\\em PLN},
        5/2022/{50 anos\\de Prolog}} {
      \fill[accent] (\x,0) circle (3.2pt);
      \node[above=5pt, font=\bfseries\small] at (\x,0) {\ano};
      \node[below=6pt, align=center, font=\scriptsize] at (\x,0) {\txt};
    }
  \end{tikzpicture}
  \vspace{12pt}

  \raggedright
  \begin{block}{Linguagem escolhida: SWI-Prolog}
    \small
    \begin{itemize}
      \item livre, bem documentado e mantido ativamente
      \item servidor HTTP, JSON e testes (\texttt{plunit}) na própria biblioteca
      \item família lógica também inclui Datalog, Mercury e Curry
    \end{itemize}
  \end{block}
\end{frame}

\begin{frame}[fragile]{Fatos, regras e consultas}
  \begin{columns}[T]
    \begin{column}{0.54\textwidth}
      \small\textbf{Base de conhecimento}
\begin{lstlisting}[style=pl]
% fatos: o que é verdade
transferencia(pedro, bruno, 500).
transferencia(bruno, carla, 300).
transferencia(bruno, davi, 15000).
conta_risco(davi).

% regras: o que se conclui
suspeita(Origem, Destino) :-
    transferencia(Origem, Destino, Valor),
    conta_risco(Destino),
    Valor > 10000.

alerta(Conta) :- suspeita(Conta, _).
\end{lstlisting}
    \end{column}
    \begin{column}{0.42\textwidth}
      \small\textbf{Consultas}
\begin{lstlisting}[style=pl]
?- suspeita(X, Y).
X = bruno, Y = davi.

?- alerta(bruno).
true.

?- alerta(pedro).
false.
\end{lstlisting}
      \vspace{6pt}
      \footnotesize\textbf{Como ler}\\[2pt]
      \begin{tabular}{@{}l l@{}}
        \texttt{:-} & ``se'' \\
        \texttt{,} & ``e'' \\
        \textcolor{accent}{\texttt{Maiúscula}} & variável \\
        \texttt{\_} & ``tanto faz''
      \end{tabular}
    \end{column}
  \end{columns}
  \fonte{\texttt{alerta/1} usa \texttt{suspeita/2}: regras se constroem sobre regras (inferência).}
\end{frame}

\begin{frame}[fragile]{Unificação: tornar dois termos iguais}
  \begin{columns}[T]
    \begin{column}{0.47\textwidth}
      \small Funciona nos dois sentidos: não há lado de ``entrada''.
\begin{lstlisting}[style=pl]
?- f(X, b) = f(a, Y).
X = a, Y = b.
\end{lstlisting}
      \vspace{8pt}
      A mesma variável duas vezes obriga os valores a serem iguais:
\begin{lstlisting}[style=pl]
autotransferencia(Conta) :-
    transferencia(Conta, Conta, _).

?- autotransferencia(C).
C = carla.
\end{lstlisting}
    \end{column}
    \begin{column}{0.47\textwidth}
      \begin{alertblock}{Por que isso é uma vantagem}
        \small
        \begin{itemize}
          \item nenhum \texttt{if}, nenhum \texttt{==}
          \item a regra \emph{só casa} quando origem e destino são iguais
          \item a mesma operação serve para comparar, extrair e construir dados
        \end{itemize}
      \end{alertblock}
      \vspace{6pt}
      {\small É o mecanismo por trás de todo o resto: unificando a pergunta com
      os fatos e com as cabeças das regras, o Prolog encontra as respostas.}
    \end{column}
  \end{columns}
\end{frame}

\begin{frame}[fragile]{Backtracking: todas as respostas, sem laço}
  \begin{columns}[T]
    \begin{column}{0.67\textwidth}
      \centering
      \begin{tikzpicture}[
          n/.style={font=\ttfamily\scriptsize, fill=codebg, rounded corners=2pt,
                    inner sep=3pt},
          e/.style={-{Stealth[length=1.8mm]}, ink!50},
          back/.style={-{Stealth[length=1.8mm]}, dashed, accent, bend right=28}]
        \node[n, fill=ink, text=paper] (r) at (3.35,0) {?- suspeita(X, Y)};
        \node[n] (a) at (0,-1.3)    {X=pedro, Y=bruno};
        \node[n] (b) at (3.35,-1.3) {X=bruno, Y=carla};
        \node[n] (c) at (6.7,-1.3)  {X=bruno, Y=davi};
        \node[n] (a2) at (0,-2.5)    {conta\_risco(bruno) \errado};
        \node[n] (b2) at (3.35,-2.5) {conta\_risco(carla) \errado};
        \node[n] (c2) at (6.7,-2.5)  {conta\_risco(davi) \certo};
        \node[n] (c3) at (6.7,-3.5)  {15000 > 10000 \certo};
        \node[n, fill=accent!15] (c4) at (6.7,-4.5) {X = bruno, Y = davi};
        \draw[e] (r) -- (a); \draw[e] (r) -- (b); \draw[e] (r) -- (c);
        \draw[e] (a) -- (a2); \draw[e] (b) -- (b2); \draw[e] (c) -- (c2);
        \draw[e] (c2) -- (c3); \draw[e] (c3) -- (c4);
        \draw[back] (a2.south) to node[below, font=\tiny]{volta} (b2.south);
        \draw[back] (b2.south) to node[below, font=\tiny]{volta} (c2.south west);
      \end{tikzpicture}
    \end{column}
    \begin{column}{0.31\textwidth}
      \small
      Falhou? O Prolog \alert{volta ao último ponto de escolha} e tenta a próxima
      alternativa.
      \vspace{8pt}

      \texttt{findall/3} recolhe todas as respostas:
\begin{lstlisting}[style=pl, style=pequeno]
?- findall(D,
     transferencia(bruno, D, _),
     L).
L = [carla, davi].
\end{lstlisting}
      \vspace{4pt}
      {\footnotesize\color{muted} A busca em profundidade com retrocesso é o
      ``laço'' do paradigma, e já vem pronta.}
    \end{column}
  \end{columns}
\end{frame}

\begin{frame}[fragile]{Recursão no lugar de \texttt{for} e \texttt{while}}
  \begin{columns}[T]
    \begin{column}{0.52\textwidth}
\begin{lstlisting}[style=pl]
% caso base: transferência direta
alcanca(A, B) :-
    transferencia(A, B, _).
% caso recursivo: passa por alguém
alcanca(A, B) :-
    transferencia(A, X, _),
    alcanca(X, B).
\end{lstlisting}
      \vspace{4pt}
\begin{lstlisting}[style=pl]
?- alcanca(pedro, davi).
true.

?- findall(B, alcanca(pedro, B), L).
L = [bruno, carla, davi].
\end{lstlisting}
      {\scriptsize\color{muted} \texttt{exemplos/05}: pedro $\to$ bruno $\to$ carla $\to$ davi}
    \end{column}
    \begin{column}{0.44\textwidth}
      \small
      Lida em voz alta, a regra é a própria definição:
      \vspace{4pt}

      \emph{``A alcança B se A transfere direto para B, \textbf{ou} se A
      transfere para algum X que alcança B.''}
      \vspace{10pt}

      \begin{itemize}
        \item não existe contador nem condição de parada escrita à mão
        \item cada cláusula é uma alternativa; o backtracking tenta as duas
        \item é uma \alert{definição}, não um laço
      \end{itemize}
    \end{column}
  \end{columns}
\end{frame}

\begin{frame}[fragile]{Relações, não funções}
  \begin{columns}[T]
    \begin{column}{0.50\textwidth}
\begin{lstlisting}[style=pl, style=pequeno]
% para quem o Pedro mandou?
?- transferencia(pedro, X, _).
X = bruno.
% quem mandou para a Carla?
?- transferencia(X, carla, _).
X = bruno.
% junte duas listas
?- append([a], [b], L).
L = [a, b].
% que listas, juntas, dão [a, b]?
?- append(X, Y, [a, b]).
X = [], Y = [a, b] ;
X = [a], Y = [b] ;
X = [a, b], Y = [].
\end{lstlisting}
    \end{column}
    \begin{column}{0.46\textwidth}
      \small
      Uma função tem entrada e saída fixas. Uma \alert{relação} responde a
      qualquer pergunta que se faça sobre ela: o que é dado e o que é
      procurado decide-se na hora da consulta.
      \vspace{8pt}

      \begin{block}{No núcleo}
        \footnotesize
        \texttt{append/3} ``ao contrário'' gera cada conta de partida da busca
        de ciclos, junto com as contas que vêm antes dela:\\[3pt]
        \texttt{append(Antes, [Inicio|\_], Contas)}
      \end{block}
    \end{column}
  \end{columns}
\end{frame}

\begin{frame}[fragile]{Mundo fechado: \texttt{false} é ``não consegui provar''}
  \begin{columns}[T]
    \begin{column}{0.48\textwidth}
      \small
      \begin{itemize}
        \item o que não pode ser provado a partir da base é tratado como falso
        \item \texttt{\textbackslash+ Objetivo} é verdade quando o objetivo
              \emph{falha} (negação por falha, Clark, 1978)
      \end{itemize}
      \vspace{4pt}
\begin{lstlisting}[style=pl]
?- conta_risco(carla).
false.     % não está na lista

% no núcleo: só avança para
% contas ainda não visitadas
\+ memberchk(Prox, Visitados)
\end{lstlisting}
    \end{column}
    \begin{column}{0.46\textwidth}
      \begin{alertblock}{Vantagem para regras de negócio}
        \small
        Se a conta não está na lista de risco, ela \emph{não é} de risco, sem
        precisar declarar isso. Categoria sem limite? A regra de limite
        simplesmente não se aplica.
      \end{alertblock}
      \vspace{4pt}
      \begin{block}{Cuidado}
        \small
        Funciona porque a base está completa: o lote traz todos os fatos.
        E \texttt{false} não diz \emph{por que} falhou.
      \end{block}
    \end{column}
  \end{columns}
\end{frame}

\section{Prolog na prática: o núcleo}

\begin{frame}[fragile]{Da frase à cláusula}
  \begin{columns}[c]
    \begin{column}{0.36\textwidth}
      \begin{block}{Enunciado}
        \small Uma transação \textbf{viola o limite da categoria} \emph{se} o
        valor dela é maior que o limite definido para a categoria.
      \end{block}
    \end{column}
    \begin{column}{0.04\textwidth}
      \centering{\Large\color{accent}$\Rightarrow$}
    \end{column}
    \begin{column}{0.56\textwidth}
\begin{lstlisting}[style=pl]
violacao(Id, limite_categoria, Motivo) :-
    transacao(Id, _, _, Valor, Categoria),
    limite_categoria(Categoria, Limite),
    Valor > Limite,
    format(atom(Motivo),
      'R$ ~w passa do limite de R$ ~w ...',
      [Valor, Limite]).
\end{lstlisting}
      {\scriptsize\color{muted}\texttt{nucleo/regras.pl} (formatação do valor em reais omitida)}
    \end{column}
  \end{columns}
  \vspace{12pt}
  \begin{itemize}\small
    \item cada linha do corpo é uma condição da frase, ligadas por ``e''
    \item categoria sem limite? \texttt{limite\_categoria/2} falha e a regra não se
          aplica: nada de \texttt{if}, \texttt{null} ou valor padrão
    \item o limite não está no código: vem do lote, que vem do banco
  \end{itemize}
\end{frame}

\begin{frame}[fragile]{Cada regra, uma cláusula independente}
  \begin{columns}[T]
    \begin{column}{0.56\textwidth}
\begin{lstlisting}[style=pl]
violacao(Id, conta_risco, Motivo) :-
    transacao(Id, _, Destino, Valor, _),
    conta_risco(Destino),
    alerta_risco(Alerta),
    Valor > Alerta, ...

% unificação: Conta aparece duas vezes
violacao(Id, autotransferencia, Motivo) :-
    transacao(Id, Conta, Conta, _, _), ...

violacao(Id, valor_invalido, Motivo) :-
    transacao(Id, _, _, Valor, _),
    Valor =< 0, ...
\end{lstlisting}
    \end{column}
    \begin{column}{0.40\textwidth}
      \small
      \begin{itemize}
        \item nenhuma regra chama outra nem depende da ordem de execução da
              outra
        \item uma transação pode quebrar várias regras: t3 aparece em
              \emph{limite} e em \emph{conta de risco}, naturalmente
        \item a ordem das cláusulas define a ordem dos achados de uma mesma
              transação
      \end{itemize}
      \vspace{4pt}
      \destaque{\footnotesize 4 regras por transação = 4 cláusulas de
      \texttt{violacao/3}}
    \end{column}
  \end{columns}
\end{frame}

\begin{frame}[fragile]{O diagnóstico inteiro}
  \begin{columns}[T]
    \begin{column}{0.63\textwidth}
\begin{lstlisting}[style=pl, style=pequeno]
diagnostico(Achados) :-
    findall(achado(transacao(Id), Regra, Motivo),
            ( transacao(Id, _, _, _, _),
              violacao(Id, Regra, Motivo) ),
            PorTransacao),
    findall(achado(ciclo(Contas, Ids), ciclo, Motivo),
            ( ciclo(Contas),
              transacoes_do_ciclo(Contas, Ids),
              motivo_ciclo(Contas, Motivo) ),
            Ciclos),
    append(PorTransacao, Ciclos, Achados).
\end{lstlisting}
    \end{column}
    \begin{column}{0.33\textwidth}
      \small
      \begin{tabular}{@{}l l@{}}
        \errado & laço escrito à mão \\[3pt]
        \errado & lista acumuladora \\[3pt]
        \errado & ordenação \\[3pt]
        \certo & backtracking percorre\\ & transações $\times$ regras \\[3pt]
        \certo & \texttt{findall/3} junta tudo \\[3pt]
        \certo & ordem = ordem do lote,\\ & depois a das cláusulas
      \end{tabular}
    \end{column}
  \end{columns}
  \vspace{8pt}
  \destaque{\small O ``validador'' é uma \emph{pergunta}: quais \texttt{(Id, Regra,
  Motivo)} tornam \texttt{violacao/3} verdade?}
\end{frame}

\begin{frame}[fragile]{Ciclos de dinheiro: busca em profundidade em duas cláusulas}
  \begin{columns}[T]
    \begin{column}{0.60\textwidth}
\begin{lstlisting}[style=pl, style=pequeno]
% chegou de volta ao início: é um ciclo
caminho(Atual, Inicio, _, [Inicio]) :-
    aresta(Atual, Inicio).
% ou segue para uma conta ainda não visitada
caminho(Atual, Inicio, Visitados, [Prox|Resto]) :-
    aresta(Atual, Prox),
    \+ memberchk(Prox, Visitados),
    caminho(Prox, Inicio, [Prox|Visitados], Resto).

% contas antes de Inicio no lote já contam como
% visitadas: cada ciclo é encontrado uma vez só
ciclo([Inicio|Resto]) :-
    contas_do_lote(Contas),
    append(Antes, [Inicio|_], Contas),
    caminho(Inicio, Inicio, [Inicio|Antes], Resto).
\end{lstlisting}
    \end{column}
    \begin{column}{0.37\textwidth}
      \centering
      \begin{tikzpicture}[scale=0.85, transform shape,
          c/.style={circle, draw=ink!60, fill=white, font=\scriptsize,
                    minimum size=1.05cm, inner sep=0pt},
          e/.style={-{Stealth[length=2mm]}, thick},
          lbl/.style={font=\tiny, fill=paper, inner sep=1pt}]
        \node[c] (p) at (0,0) {pedro};
        \node[c] (b) at (2.6,0) {bruno};
        \node[c] (c) at (1.3,-2.0) {carla};
        \node[c, diamond, minimum size=1.1cm, draw=bad, font=\tiny, align=center]
          (f) at (4.0,-1.4) {conta\\fantasma};
        \node[c] (d) at (2.9,-3.0) {davi};
        \draw[e, accent] (p) -- node[lbl, above]{t1} (b);
        \draw[e, accent] (b) -- node[lbl, right]{t4} (c);
        \draw[e, accent] (c) -- node[lbl, left]{t5} (p);
        \draw[e, ink!45] (b) -- node[lbl, above right]{t3} (f);
        \draw[e, ink!30, dotted] (c) -- node[lbl, below left]{t6: R\$ 0} (d);
        \draw[e, ink!30, dotted] (b) to[out=60, in=120, looseness=6]
          node[lbl, above]{t2} (b);
      \end{tikzpicture}

      \footnotesize\raggedright
      \begin{itemize}
        \item a pilha da busca é o backtracking: ao falhar, o caminho se desfaz
              sozinho
        \item \texttt{Visitados} garante que a busca termina
      \end{itemize}
    \end{column}
  \end{columns}
\end{frame}

\begin{frame}[fragile]{Saída real do núcleo}
\begin{lstlisting}[style=terminal]
$ swipl nucleo/regras.pl nucleo/lotes/seed.pl
?- mostrar_diagnostico.
t1 [limite_categoria] R$ 60.000 passa do limite de R$ 50.000 da categoria suprimentos
t2 [autotransferencia] origem e destino são a mesma conta (bruno)
t3 [limite_categoria] R$ 15.000 passa do limite de R$ 10.000 da categoria servicos
t3 [conta_risco] valor alto (R$ 15.000) enviado para a conta de risco conta_fantasma
t6 [valor_invalido] valor zerado ou negativo (R$ 0)
t1,t4,t5 [ciclo] o dinheiro sai de pedro e volta para pedro (pedro → bruno → carla → pedro)
\end{lstlisting}
  \vspace{6pt}
  \begin{columns}[T]
    \begin{column}{0.55\textwidth}
      \begin{alertblock}{O núcleo achou o que nós deixamos passar}
        \small Na proposta, escrevemos à mão que t3 quebrava uma regra. Quebra
        duas: limite \textbf{e} conta de risco.
      \end{alertblock}
    \end{column}
    \begin{column}{0.41\textwidth}
      \footnotesize
      \begin{itemize}
        \item a tela mostra exatamente isso: mesmos fatos, mesmo \texttt{diagnostico/1}
        \item 50 testes \texttt{plunit} passam em 0,04~s
        \item o núcleo roda sozinho, sem a web
      \end{itemize}
    \end{column}
  \end{columns}
\end{frame}

\begin{frame}[fragile]{Uma regra nova é uma cláusula a mais}
  \begin{columns}[T]
    \begin{column}{0.50\textwidth}
\begin{lstlisting}[style=pl]
violacao(Id, nova_regra, Motivo) :-
    transacao(Id, ...),
    condicao_da_regra,
    format(atom(Motivo), ...).
\end{lstlisting}
      \vspace{6pt}
      \small
      \begin{tabular}{@{}l l@{}}
        \toprule
        \textbf{Parte} & \textbf{O que muda} \\
        \midrule
        \texttt{regras.pl} & +1 cláusula \\
        \texttt{diagnostico/1} & nada \\
        servidor HTTP / JSON & nada \\
        banco & nada \\
        BFF e tela & nome da regra \\
        \bottomrule
      \end{tabular}
    \end{column}
    \begin{column}{0.45\textwidth}
      \begin{block}{Configuração também é fato}
        \small
        Limite de \emph{serviços} de R\$ 10.000 para R\$ 100: o lote é remontado
        do banco como fatos e o diagnóstico muda.\\[4pt]
        \textbf{Nenhuma linha de código nova.}
      \end{block}
      \vspace{6pt}
      {\small O \texttt{findall/3} do diagnóstico pergunta por \emph{qualquer}
      cláusula de \texttt{violacao/3}; uma regra nova entra nele sozinha.}
    \end{column}
  \end{columns}
\end{frame}

\section{Comparação com outras linguagens}

\begin{frame}[fragile]{A mesma regra em Java (POO)}
  \begin{columns}[T]
    \begin{column}{0.55\textwidth}
\begin{lstlisting}[language=Java, style=pequeno]
interface Regra {
  Optional<Achado> verificar(Transacao t);
}

class Autotransferencia implements Regra {
  public Optional<Achado> verificar(Transacao t) {
    if (t.origem().equals(t.destino()))
      return Optional.of(
          new Achado(t.id(), "autotransferencia"));
    return Optional.empty();
  }
}

// o validador percorre tudo à mão
List<Achado> achados = new ArrayList<>();
for (Transacao t : lote)
  for (Regra r : regras)
    r.verificar(t).ifPresent(achados::add);
\end{lstlisting}
    \end{column}
    \begin{column}{0.44\textwidth}
\begin{lstlisting}[style=pl, style=pequeno]
violacao(Id, autotransferencia, _) :-
    transacao(Id, Conta, Conta, _, _).

% o "validador"
?- findall(Id-Regra,
           violacao(Id, Regra, _),
           Achados).
\end{lstlisting}
      \vspace{4pt}
      \footnotesize
      \begin{itemize}
        \item \textbf{POO:} a regra vira classe; laço, acumulador e comparação
              são escritos à mão
        \item \textbf{Prolog:} a regra vira cláusula; a comparação é a
              unificação e o laço é o backtracking
        \item \textbf{A favor da POO:} tipos estáticos, encapsulamento e
              isolamento natural (um objeto por lote)
      \end{itemize}
    \end{column}
  \end{columns}
\end{frame}

\begin{frame}[fragile]{Busca de ciclos em Python (imperativo)}
  \begin{columns}[T]
    \begin{column}{0.55\textwidth}
\begin{lstlisting}[language=Python, style=pequeno]
def ciclos(arestas, contas):
  achados = []
  for i, inicio in enumerate(contas):
    pilha = [(inicio, [inicio])]
    while pilha:
      atual, cam = pilha.pop()
      for prox in arestas.get(atual, []):
        if prox == inicio:
          achados.append(cam + [inicio])
        elif prox not in cam + contas[:i]:
          pilha.append((prox, cam + [prox]))
  return achados
\end{lstlisting}
    \end{column}
    \begin{column}{0.44\textwidth}
\begin{lstlisting}[style=pl, style=pequeno]
caminho(Atual, Inicio, _, [Inicio]) :-
    aresta(Atual, Inicio).
caminho(Atual, Inicio, Visitados,
        [Prox|Resto]) :-
    aresta(Atual, Prox),
    \+ memberchk(Prox, Visitados),
    caminho(Prox, Inicio,
            [Prox|Visitados], Resto).
\end{lstlisting}
      \vspace{6pt}
      \footnotesize
      \begin{itemize}
        \item \textbf{Python:} pilha, cópia do caminho, laço externo e lista de
              resultados, tudo controlado à mão
        \item \textbf{Prolog:} só a definição de ciclo; a pilha e o
              ``desfazer'' vêm do backtracking
        \item não é só tamanho: é \alert{quem} cuida da busca
      \end{itemize}
    \end{column}
  \end{columns}
\end{frame}

\begin{frame}[fragile]{SQL: o primo declarativo}
  \begin{columns}[T]
    \begin{column}{0.52\textwidth}
\begin{lstlisting}[language=SQL, style=pequeno, morekeywords={recursive}]
with recursive alcanca(origem, destino) as (
    select origem, destino from transacoes
  union
    select a.origem, t.destino
    from alcanca a
    join transacoes t on t.origem = a.destino
)
select destino from alcanca
where origem = 'pedro';
\end{lstlisting}
      \vspace{4pt}
\begin{lstlisting}[style=pl, style=pequeno]
alcanca(A, B) :- transferencia(A, B, _).
alcanca(A, B) :- transferencia(A, X, _),
                 alcanca(X, B).

?- alcanca(pedro, B).
\end{lstlisting}
    \end{column}
    \begin{column}{0.45\textwidth}
      \small
      \begin{itemize}
        \item SQL também diz \emph{o quê}, não \emph{como}, e também vem da
              lógica (cálculo relacional); Datalog liga os dois mundos
        \item em SQL a recursão é uma construção especial; em Prolog é o jeito
              normal de escrever, com listas e termos como
              \texttt{ciclo([pedro, \ldots])}
        \item \textbf{A favor do SQL:} volume, persistência e concorrência.
              Por isso os dados moram no Postgres
      \end{itemize}
    \end{column}
  \end{columns}
\end{frame}

\begin{frame}{Comparação geral}
  \footnotesize
  \renewcommand{\arraystretch}{1.3}
  \begin{tabularx}{\linewidth}{@{}>{\bfseries}l P Y Y@{}}
    \toprule
    & \textbf{Prolog} & \textbf{Java (POO)} & \textbf{SQL} \\
    \midrule
    Você descreve & fatos e regras & objetos e métodos & consultas sobre tabelas \\
    Repetição & recursão e backtracking & laços & \texttt{WITH RECURSIVE} \\
    Busca de soluções & nativa (resolução) & escrita à mão & só sobre os dados \\
    Sentido & relação: vários sentidos & entrada $\to$ saída & tabelas $\to$ tabela \\
    Tipos & dinâmicos & estáticos & esquema das tabelas \\
    Brilha em & regras, grafos, busca & sistemas grandes com estado & volume e persistência \\
    \bottomrule
  \end{tabularx}
  \normalsize
  \destaque{\small Nenhuma ganha em tudo: Prolog ganha quando o problema \emph{é feito de regras}.}
\end{frame}

\section{Dificuldades, aplicações e conclusões}

\begin{frame}{Dificuldades}
  \begin{columns}[T]
    \begin{column}{0.48\textwidth}
      \small
      \textbf{\color{accent}Fatos são globais}\\
      isolar cada lote exigiu \texttt{assert}/\texttt{retract}, mutex e
      \texttt{setup\_call\_cleanup}: um lote por vez
      \vspace{7pt}

      \textbf{\color{accent}Terminação é responsabilidade sua}\\
      \texttt{alcanca/2} em grafo com ciclo não para; foi preciso a lista de
      visitadas
      \vspace{7pt}

      \textbf{\color{accent}A busca pode explodir}\\
      listar todos os ciclos é exponencial: 40 contas e 100 transações levavam
      20~s; com poda, 2~s (cerca de 14.700 ciclos)
    \end{column}
    \begin{column}{0.48\textwidth}
      \small
      \textbf{\color{accent}Sem tipos}\\
      erro de digitação ou biblioteca não carregada vira \emph{``Unknown
      procedure''} só ao rodar os testes
      \vspace{7pt}

      \textbf{\color{accent}\texttt{false} não explica}\\
      a resposta não diz onde a prova falhou; depurar exige \texttt{trace}
      \vspace{7pt}

      \textbf{\color{accent}Nem tudo é lógica}\\
      ler e conferir o JSON de entrada ficou procedural, campo a campo
    \end{column}
  \end{columns}
\end{frame}

\begin{frame}[fragile]{Tabling: quando a linguagem devolve o controle}
  \begin{columns}[T]
    \begin{column}{0.48\textwidth}
      \small\textbf{Grafo com ciclo, sem tabling}
\begin{lstlisting}[style=pl]
transferencia(pedro, bruno, 500).
transferencia(bruno, carla, 300).
transferencia(carla, pedro, 200).

?- findall(B, alcanca(pedro, B), L).
% não termina: pedro, bruno, carla,
% pedro, bruno... para sempre
\end{lstlisting}
    \end{column}
    \begin{column}{0.48\textwidth}
      \small\textbf{Com uma diretiva a mais}
\begin{lstlisting}[style=pl]
:- table alcanca/2.

?- findall(B, alcanca(pedro, B), L).
L = [pedro, carla, bruno].
\end{lstlisting}
      \vspace{4pt}
      \footnotesize
      O SWI-Prolog memoriza as respostas já encontradas (resolução SLG): ciclos e
      recursão à esquerda passam a terminar, como o \texttt{UNION} do SQL.
    \end{column}
  \end{columns}
  \vspace{8pt}
  \begin{block}{Por que o núcleo não usa}
    \small O diagnóstico precisa do \emph{caminho} de cada ciclo, não só de quem
    alcança quem; por isso a busca guarda as contas visitadas.
  \end{block}
\end{frame}

\begin{frame}{Onde o paradigma brilha e onde não}
  \begin{columns}[T]
    \begin{column}{0.50\textwidth}
      \begin{alertblock}{Brilha}
        \small
        \begin{itemize}
          \item regras de negócio e legislação\\
                {\scriptsize\color{muted} lei britânica de nacionalidade como programa (1986)}
          \item sistemas especialistas e diagnóstico
          \item grafos: caminhos, dependências, ciclos
          \item planejamento e escalonamento
          \item linguagem natural\\
                {\scriptsize\color{muted} análise de perguntas no IBM Watson (2011)}
        \end{itemize}
      \end{alertblock}
    \end{column}
    \begin{column}{0.44\textwidth}
      \begin{block}{Não é a melhor escolha}
        \small
        \begin{itemize}
          \item interfaces gráficas
          \item cálculo numérico pesado
          \item programas com muito estado mutável
          \item grande volume de dados persistidos
        \end{itemize}
      \end{block}
    \end{column}
  \end{columns}
  \vspace{8pt}
  \destaque{\small No nosso sistema: \textbf{React} na tela, \textbf{Postgres} nos
  dados, \textbf{Prolog} nas regras.}
\end{frame}

\begin{frame}{Respondendo à pergunta do trabalho}
  \centering
  {\small\emph{Como o paradigma influencia a maneira de pensar, modelar e
  implementar uma solução?}}
  \vspace{10pt}

  \begin{columns}[T]
    \begin{column}{0.31\textwidth}
      \begin{alertblock}{Pensar}
        \small De ``como calcular?'' para \textbf{``o que é verdade?''}. O
        enunciado do problema já é quase a especificação.
      \end{alertblock}
    \end{column}
    \begin{column}{0.31\textwidth}
      \begin{alertblock}{Modelar}
        \small \textbf{Fatos e relações} em vez de objetos:
        \texttt{transacao/5}, \texttt{aresta/2}, \texttt{ciclo/1}. Os dados já
        são lógica.
      \end{alertblock}
    \end{column}
    \begin{column}{0.31\textwidth}
      \begin{alertblock}{Implementar}
        \small \textbf{Cláusulas + busca automática}. O controle volta onde
        importa: terminação e poda.
      \end{alertblock}
    \end{column}
  \end{columns}
\end{frame}

\begin{frame}{Conclusão}
  \begin{itemize}
    \setlength{\itemsep}{7pt}
    \item As regras ficaram \textbf{parecidas com o enunciado} e independentes entre si
    \item \textbf{Unificação e backtracking} substituíram comparações, laços e acumuladores
    \item A busca de ciclos, que em Java ou Python seria uma DFS escrita à mão,
          virou \textbf{duas cláusulas}
    \item O núcleo encontrou um achado que a especificação feita à mão deixou passar
    \item O preço: estado global, terminação e desempenho continuam sendo
          \textbf{responsabilidade nossa}
  \end{itemize}
  \vspace{10pt}
  \destaque{Prolog não substitui as outras linguagens: é a ferramenta certa para
  a parte do problema que é \textbf{feita de regras}.}
\end{frame}

\begin{frame}[standout]
  Obrigado!\\[10pt]
  {\normalsize\ttfamily ?- pergunta(X).}\\[18pt]
  {\scriptsize\normalfont github.com/ThiaagoMP/trabalho-final-linguagens}
\end{frame}

\appendix

\begin{frame}{Referências}
  \scriptsize
  \begin{itemize}
    \setlength{\itemsep}{2pt}
    \item CLARK, K. L. Negation as failure. In: \emph{Logic and Data Bases}. Plenum, 1978.
    \item COLMERAUER, A.; ROUSSEL, P. The birth of Prolog. \emph{ACM SIGPLAN Notices}, 28(3), 1993.
    \item KÖRNER, P. et al. Fifty years of Prolog and beyond. \emph{TPLP}, 22(6), 2022.
    \item KOWALSKI, R. Predicate logic as programming language. \emph{IFIP Congress}, 1974.
    \item KOWALSKI, R. Algorithm = logic + control. \emph{Communications of the ACM}, 22(7), 1979.
    \item LALLY, A.; FODOR, P. Natural language processing with Prolog in the IBM Watson system. \emph{ALP Newsletter}, 2011.
    \item ROBINSON, J. A. A machine-oriented logic based on the resolution principle. \emph{Journal of the ACM}, 12(1), 1965.
    \item SERGOT, M. J. et al. The British Nationality Act as a logic program. \emph{Communications of the ACM}, 29(5), 1986.
    \item WIELEMAKER, J. et al. SWI-Prolog. \emph{TPLP}, 12(1--2), 2012.
    \item SWI-Prolog Reference Manual, seção \emph{Tabled execution}. \url{https://www.swi-prolog.org}
  \end{itemize}
\end{frame}

\end{document}
